% Session 7 -- what an MPS of bond dimension M can and cannot represent.
% Placement: sec. "Why MPS work in 1D, and where they stop working", after (7.10).
%
% LAYOUT NOTE: the only pgfplots figure of this set. Everything, annotations
% included, is positioned in (axis cs:...) so nothing can drift off the curves.
\begin{aqmfigure}{fig:s07:ceiling}{The bond dimension is an upper limit on the entanglement, and this
determines when matrix product states work. An MPS with bond dimension $M$ has at most $M$
Schmidt terms across any bond, so by (7.10) it can have at most $S=\log M$ there, which is shown by
the two horizontal lines. A gapped one-dimensional ground state follows the area law of Section 7.6,
and its entropy saturates at a value set by the correlation length and not by the subsystem size, so
it stays below a fixed and moderate limit for any chain length. A critical chain violates the area law
logarithmically, (7.12), so it goes above any fixed limit at some point, and $M$ must be increased as
the system grows, slowly but without bound. A general state of the Hilbert space of dimension $2^L$
follows a volume law, $S=\ell\log 2$, and it goes above any usable $M$ after a few sites. This is why
time evolution over long times, highly excited states and two-dimensional systems are all out of
reach. The entropy is drawn in bits, $S/\log 2$, with central charge $c=1$, and the curves are
schematic and not fits to a specific model.}

\begin{tikzpicture}
\begin{axis}[
  aqmplot,
  width=12.4cm, height=7.0cm,
  axis lines=left,
  axis line style={aqm, black},
  xlabel={subsystem size $\ell$ (sites)},
  ylabel={$S/\log 2$},
  xlabel style={at={(axis description cs:0.5,-0.11)}, anchor=north},
  ylabel style={at={(axis description cs:-0.075,0.5)}, anchor=south, rotate=90},
  xmin=0, xmax=74, ymin=0, ymax=6.5,
  xtick={0,16,32,48,64}, ytick={0,2,4,6},
  enlargelimits=false,
]

% --- the two MPS ceilings, S = log M ----------------------------------------
\addplot[aqmguide, forget plot, domain=0:74, samples=2] {2};
\addplot[aqmguide, forget plot, domain=0:74, samples=2] {4};
\node[aqmlabsmall, anchor=south west, text=aqmmuted] at (axis cs:6.5,2.04)
  {$\log M$ for $M=4$};
\node[aqmlabsmall, anchor=south west, text=aqmmuted] at (axis cs:6.5,4.04)
  {MPS ceiling $S=\log M$, $\;M=16$};

% --- volume law: a generic state --------------------------------------------
\addplot[aqmcurvebroken, dash pattern=on 4pt off 2pt, domain=0:6.4, samples=2] {x};
\node[aqmlabsmall, anchor=west, align=left, text=aqmbroken] at (axis cs:7.6,5.62)
  {generic (volume-law) state: $S=\ell\log 2$\\[-1pt]
   no practical $M$ survives past a few sites};

% --- critical 1D chain: (c/3) log l ------------------------------------------
\addplot[black, smooth, dash pattern=on 2.5pt off 2pt, domain=1:74, samples=120]
  {0.4809*ln(x)+0.35};
\node[aqmlabsmall, anchor=north east, align=right] at (axis cs:73,1.92)
  {critical 1D chain, $S=\frac{c}{3}\log\ell$ (7.12):\\[-1pt]
   grows without bound, so $M$ must grow too};

% --- gapped 1D ground state: saturates ---------------------------------------
\addplot[aqmcurve, domain=0:74, samples=140] {0.90*(1-exp(-x/3.2))};
\node[aqmlabsmall, anchor=north west, align=left, text=aqmaccent] at (axis cs:26,0.71)
  {gapped 1D ground state: $S$ saturates};

% --- where the small ceiling is exhausted ------------------------------------
\node[circle, draw=black, fill=white, line width=\aqmthinwidth, minimum size=5pt,
      inner sep=0pt] (xcross) at (axis cs:30.9,2) {};
\draw[aqmguide, draw=aqmmuted] (axis cs:30.9,2.78) -- (xcross);
\node[aqmlabsmall, anchor=north west, align=left] at (axis cs:19,3.52)
  {$M=4$ is exhausted here; a larger $M$ lifts the\\[-1pt]
   ceiling, but it never flattens the curve};

\end{axis}
\end{tikzpicture}

\end{aqmfigure}
